\begin{center}
    \textbf{RESUMO}
\end{center}

\noindent Em parceria com a Riachuelo, a equipe desenvolveu o conceito de uma etiqueta antifurto que o próprio cliente libera depois de pagar pelo aplicativo, para reduzir as filas no varejo de moda. A definição do mercado delimitou como segmento-alvo as redes de moda \textit{mass market} que já utilizam self-checkout, e a expansão das lojas conceito da Riachuelo prevista para 2026 foi identificada como o mercado obtível no curto prazo. A pesquisa com consumidores combinou um survey com 142 respostas e onze entrevistas semiestruturadas. A clusterização exploratória identificou quatro grupos: dependentes de atendimento (39\%), autônomos digitais (21\%), frustrados com disponibilidade (9\%) e pragmáticos equilibrados (31\%). Apenas 24\% dos respondentes endossaram a compra sem checagem manual na saída, com aceitação total entre os autônomos digitais e rejeição quase unânime entre os dependentes de atendimento. Os achados foram traduzidos em oito necessidades, sete desejos e cinco demandas, e sete dessas vozes foram desdobradas em requisitos técnicos por meio da casa da qualidade, com benchmarking de Renner, C\&A e Zara, no qual a liberação da peça condicionada ao pagamento, o tempo de pagamento e a clareza do fluxo no aplicativo concentraram cerca de 71\% da importância técnica. A equipe também elaborou o esquema à mão livre de um dispositivo de 98 por 82 por 26 mm, com retenção por embreagem de esferas e liberação autorizada pelo celular. A equipe concluiu que a solução tende a gerar valor quando preserva a opção de atendimento humano.

\vspace{1em}
\noindent\textbf{Palavras-chave:} varejo de moda; etiqueta antifurto; autoatendimento.
